\section{An auxiliary prescribed-volume solution}
\label{sec:auxiliary-volume}

We retain the fibration and prepared forms supplied by
\Cref{prop:regular-preparation}.  Thus the closed definite pair $(A,C)$ has
the original periods, both forms vanish on the fibers, and on the regular
locus they are invariant under the fiber torus translations.  The closed
invariant form $\Psi$ has fiber integral one and satisfies
\begin{equation}
  \Psi\wedge A=\Psi\wedge C=0.
  \label{eq:auxiliary-prepared-orthogonality}
\end{equation}
Near each punctured nodal fiber it is the form $\dd a\wedge\dd b$ of
\Cref{thm:nodal-model}.  All the model disks, smaller disks, and annuli in
this section are fixed before the parameter $\eps$ is chosen.  We write
$F$ for the Poincar\'e dual of an oriented fiber and
\[
  V=[A]^\perp\cap[C]^\perp.
\]
In particular $F\in V$, $F^2=0$, and $[B]F>0$.

\begin{theorem}[Auxiliary volume equation]
\label{thm:auxiliary-solution}
For every sufficiently small $\eps>0$ there is a smooth closed real
two-form $D_\eps$ such that
\begin{equation}
  A\wedge D_\eps=C\wedge D_\eps=0,
  \qquad D_\eps^2=A^2,
  \label{eq:auxiliary-exact-equations}
\end{equation}
and $(A,C,D_\eps)$ is a positive triple.  Its sign agrees with positive
area on the oriented fibers.  The construction uses an auxiliary
cohomology class, whose relation to the required class $[B]$ is stated in
\Cref{prop:auxiliary-class}.
\end{theorem}

The form $D_\eps$ supplies a known taming class for the later pair
deformation.  To recover the original third class, we need the following
comparison, with one choice of $\eps$ valid for every curve that may
appear along that deformation.

\begin{proposition}[Class and chamber of the auxiliary form]
\label{prop:auxiliary-class}
Fix a norm on the finite-dimensional space $V$.  The forms of
\Cref{thm:auxiliary-solution} can be chosen so that
\begin{equation}
\begin{split}
  [D_\eps]&=\lambda_\eps F+\eps u_\eps,\qquad u_\eps\in V,\\
  \|u_\eps\|&\le K,\qquad u_\eps F\ge a_0>0,\\
  c_0\eps^{-1}&\le\lambda_\eps\le c_1\eps^{-1},
\end{split}
\label{eq:auxiliary-class-expansion}
\end{equation}
for fixed positive constants and all sufficiently small $\eps$.
Moreover $[D_\eps]$ and $[B]$ belong to the same component of the
positive cone of $V$.

There is a single sufficiently small choice of $\eps$ with the
following property.  For every definite closed pair with periods
$[A],[C]$, with either associated almost-complex sign, and every
irreducible closed curve for that structure with class
$d\in H^2(X;\Z)/\mathrm{torsion}$,
\begin{equation}
  [D_\eps]d>0\quad\Longrightarrow\quad[B]d>0.
  \label{eq:auxiliary-uniform-transfer}
\end{equation}
The choice depends only on the prepared construction and its period
data, and is independent of the later pair or curve.
\end{proposition}

We first glue an invariant solution on the regular region to the exact
nodal models and impose three cohomological moments.  The remaining
errors are exponential in $\eps^{-1}$, whereas the inverse and
multiplication estimates needed to correct them have only polynomial
losses.  This balance between exponential gluing error and polynomial
analytic loss follows the collapsing strategy of Gross--Wilson
\cite[Sections~4--5]{GrossWilson2000}.  Here the correction is an exact
two-form: a global complex structure, needed for a scalar complex
Monge--Amp\`ere equation, is not yet available.  The final subsection
proves \Cref{prop:auxiliary-class} by computing $[D_\eps]$ from the
prepared fibration and excluding all curve classes with negative fiber
pairing at once.

\subsection{The invariant solution and its gluing}

On the regular part define
\begin{equation}
\begin{split}
  \beta_\eps
    &=\frac{1}{2\eps}p_*(A^2)
      -\frac{\eps}{2}p_*(\Psi^2),\\
  D^{\mathrm{sf}}_\eps&=\eps\Psi+p^*\beta_\eps.
\end{split}
\label{eq:auxiliary-semiflat}
\end{equation}
Here $p_*$ is integration along the oriented torus fibers.  Every
two-form on the two-dimensional base is closed, so
$D^{\mathrm{sf}}_\eps$ is closed.  Since $A$ and $C$ have zero vertical
restriction, their products with a basic two-form vanish.  Thus
\Cref{eq:auxiliary-prepared-orthogonality} gives
\[
  D^{\mathrm{sf}}_\eps\wedge A
  =D^{\mathrm{sf}}_\eps\wedge C=0.
\]
For any invariant top form $\zeta$ on the regular part,
\begin{equation}
  \zeta=\Psi\wedge p^*(p_*\zeta).
  \label{eq:auxiliary-fiber-integration}
\end{equation}
This is the invariant fiber-integration identity
\eqref{eq:regular-invariant-top-form} for the prepared form $\Psi$.
Expanding the square in \Cref{eq:auxiliary-semiflat} and using
\Cref{eq:auxiliary-fiber-integration}, we obtain
\begin{equation}
\begin{split}
  (D^{\mathrm{sf}}_\eps)^2
   &=\eps^2\Psi^2+2\eps\Psi\wedge p^*\beta_\eps\\
   &=\eps^2\Psi^2+
      \Psi\wedge p^*\bigl(p_*(A^2)-\eps^2p_*(\Psi^2)\bigr)
     =A^2.
\end{split}
\label{eq:auxiliary-semiflat-square}
\end{equation}
The resulting triple is positive, and the sign of its third form is
specified by its vertical restriction, which has integral $\eps>0$.

In a nodal disk with holomorphic volume form
$\Sigma=h\,\dd s\wedge\dd z$, the prepared data are
\[
  A=\operatorname{Re}\Sigma,\qquad
  C=x_0\operatorname{Re}\Sigma+y_0\operatorname{Im}\Sigma,
  \qquad y_0\ne0,
\]
with $x_0,y_0$ constant.  On its regular part
$\Psi=\dd a\wedge\dd b$, so $\Psi^2=0$.  Consequently
\Cref{eq:auxiliary-semiflat} is precisely
\begin{equation}
  \eps\,\dd a\wedge\dd b
  +\frac{|h|^2}{\eps}\operatorname{Im}\tau(s)\,
       \dd x_1\wedge\dd x_2,
  \qquad s=x_1+i x_2.
  \label{eq:auxiliary-model-semiflat}
\end{equation}
The form $D^{\mathrm{loc}}_\eps$ of \Cref{thm:nodal-model} has the same
two periods on a regular annular torus bundle and differs from
\Cref{eq:auxiliary-model-semiflat} exponentially in every fixed finite
number of derivatives.

Here is an explicit bounded primitive construction for this gluing.
After choosing a radial trivialization, a fixed annular bundle is
$M\times[r_0,r_1]$, where $M$ is its compact three-dimensional mapping
torus.  If $\zeta$ is an exact closed two-form on this annulus, radial
homotopy gives
\[
  \zeta=\pi_M^*(\zeta|_{M\times\{r_*\}})+\dd K\zeta,
  \qquad
  K\zeta(r)=\int_{r_*}^{r}\iota_{\partial_r}\zeta(u)\,\dd u.
\]
The restriction at $r_*$ is exact.  For a fixed metric on $M$, let
$G_M$ be the Green operator on the complement of harmonic forms.
Then $\dd_M^*G_M(\zeta|_{r_*})$ is a primitive of that restriction.
Adding its pullback to $K\zeta$ produces a primitive on a slightly
smaller annulus.  Fixed-geometry elliptic estimates on $M$ and the radial
integral bound its $C^j$ norm by finitely many $C^\ell$ norms of $\zeta$.
All these operators and domains are independent of $\eps$.

Apply this construction to
$D^{\mathrm{loc}}_\eps-D^{\mathrm{sf}}_\eps$.  Exactness follows from
the two annular period equalities in \Cref{thm:nodal-model}.  We obtain
a primitive $\alpha_\eps$ with exponential finite-order bounds.  For a
fixed cutoff $\chi$ equal to one at the inner edge and zero at the outer
edge, replace $D^{\mathrm{sf}}_\eps$ on the annulus by
\[
  D^{\mathrm{sf}}_\eps+\dd(\chi\alpha_\eps).
\]
It equals $D^{\mathrm{loc}}_\eps$ at the inner edge and
$D^{\mathrm{sf}}_\eps$ at the outer edge.  Filling the inner disk with
$D^{\mathrm{loc}}_\eps$ gives a smooth closed form $E_\eps$ on $X$.
The differences from the exact equations are supported on the fixed
gluing annuli and have exponential finite-order bounds there.  In
particular the wedge Gram matrix remains positive for small $\eps$.
When there are no nodal fibers, simply take
$E_\eps=D^{\mathrm{sf}}_\eps$ on all of $X$.

The remaining correction will be exact.  It therefore cannot change
$\int_X A\wedge D$, $\int_X C\wedge D$, or $\int_X D^2$ for a closed
form $D$.  We must first impose the values required by the three target
equations: zero for the first two integrals and one for the third.  Set
\begin{equation}
\begin{aligned}
  a_\eps&=\int_X E_\eps\wedge A,&
  b_\eps&=\int_X E_\eps\wedge C,\\
  \widehat E_\eps&=E_\eps-a_\eps A-b_\eps C,&
  q_\eps&=\int_X\widehat E_\eps^2,\\
  r_\eps&=q_\eps^{-1/2},&
  D^0_\eps&=r_\eps\widehat E_\eps .
\end{aligned}
\label{eq:auxiliary-normalization}
\end{equation}
The period normalization gives
$\int A^2=\int C^2=1$ and $\int A\wedge C=0$.  Thus
\begin{equation}
  \int_X A\wedge D^0_\eps
  =\int_X C\wedge D^0_\eps=0,\qquad
  \int_X(D^0_\eps)^2=1.
  \label{eq:auxiliary-exact-moments}
\end{equation}
The approximate equations give
$a_\eps,b_\eps=O(e^{-c/\eps})$ and
$q_\eps=1+O(e^{-c/\eps})$, after reducing $c>0$ if necessary.
Hence $q_\eps>0$ and $r_\eps=1+O(e^{-c/\eps})$.
Subtracting constant multiples of $A,C$ and multiplying the third form
by a positive constant do not change the pointwise three-plane spanned
by the triple.

\subsection{Polynomial geometry and Sobolev bounds}

Let $g_\eps$ be the metric whose self-dual three-plane is
$\operatorname{span}(A,C,D^0_\eps)$ and whose volume form is
\begin{equation}
  \nu=\frac{A^2}{2}.
  \label{eq:auxiliary-fixed-volume}
\end{equation}
The preceding span observation allows us to use $E_\eps$ in this
definition as well.  In particular $\vol_{g_\eps}(X)=1/2$.
Sobolev norms in the rest of this section are intrinsic norms for
$g_\eps$, using its covariant derivative and volume.

\begin{lemma}[Bounds for the approximate triple]
\label{lem:auxiliary-geometry}
For every fixed nonnegative integer $k$, the following statements hold
for sufficiently small $\eps$.
\begin{enumerate}[label=\textup{(\roman*)}]
\item
The metric has a finite coordinate cover whose cardinality, inverse
coordinate radii, inverse interior margins, metric and inverse-metric
comparisons, and coefficient derivatives through any prescribed finite
order are bounded by powers of $\eps^{-1}$.
\item
The frame $(A,C,D^0_\eps)$ and the inverse of its pointwise Gram matrix
have polynomial intrinsic derivative bounds through order $k$.  The
Gram matrix itself is uniformly positive and bounded.
\item
Sobolev embedding and multiplication at a fixed sufficiently large
order have polynomial constants.  In particular, for an integer
$k\ge6$ there are constants $C,q$ such that
\begin{equation}
  \|h\wedge \widetilde h\|_{H^k}
   \le C\eps^{-q}\|h\|_{H^k}\|\widetilde h\|_{H^k}
  \label{eq:auxiliary-product}
\end{equation}
for two-forms $h,\widetilde h$.
\item
The three free residuals
\begin{equation}
  \mathcal R_\eps=
   \left(-A\wedge D^0_\eps,\,
         -C\wedge D^0_\eps,\,
         \frac{A^2-(D^0_\eps)^2}{2}\right)
  \label{eq:auxiliary-residual}
\end{equation}
satisfy, for some $C,s,c>0$ depending on the chosen finite order,
\begin{equation}
  \|\mathcal R_\eps\|_{H^k}
      \le C\eps^{-s}e^{-c/\eps}.
  \label{eq:auxiliary-residual-estimate}
\end{equation}
\end{enumerate}
\end{lemma}

\begin{proof}
We first check the three types of region in the construction.  On a
model core, $A$ and $C$ are fixed constant linear combinations of the
two parallel real components of $\Sigma$, and $E_\eps$ is the local
parallel K\"ahler form.  The model metric has volume $A^2/2$.
Consequently it is exactly $g_\eps$ there.  The intrinsic charts and
bounds of \Cref{thm:nodal-model} therefore apply.  This argument
transports the metric, forms, and charts together through the
holomorphic identification of the model disks.  It requires no bound
on derivatives of that identification in a separate fixed smooth
atlas.

On the complement of smaller nodal disks, use finitely many fixed
regular bundle charts.  The coefficients in
\Cref{eq:auxiliary-semiflat} and all their fixed-order derivatives
grow at most polynomially in $\eps^{-1}$.  Its products with $A,C$
vanish and its square is $A^2$.  The remaining two-by-two wedge block
is that of the fixed definite pair.  It is uniformly positive on this
compact region, and the same assertion holds on the cores because
there its coefficients are fixed constants.  The exponential annular
gluing errors preserve these bounds on all of $X$.

For completeness, the passage from these form bounds to metric bounds
is algebraic.  Normalize the frame by the positive square root of its
wedge Gram matrix, obtaining $\Theta_1,\Theta_2,\Theta_3$ with
$\Theta_i\wedge\Theta_j=2\delta_{ij}\nu$ and with the frame orientation
chosen for a positive metric.  In oriented coordinates, the identity
\begin{equation}
  g_\eps(v,w)\nu
   =\frac16\sum_{i,j,\ell=1}^3
       \epsilon_{ij\ell}
       (\iota_v\Theta_i)\wedge(\iota_w\Theta_j)\wedge\Theta_\ell
  \label{eq:auxiliary-metric-algebra}
\end{equation}
follows by evaluating on a standard self-dual orthonormal frame.
In the fixed regular charts this gives polynomial coefficient and
derivative bounds for $g_\eps$.  Its determinant in these charts is
fixed by $\nu$, so the adjugate formula gives polynomial bounds for
the inverse metric as well.  The core metric and inverse bounds were
already supplied by the transported model charts.  Matrix square
roots and inverses have bounded derivatives on the uniformly positive
Gram matrices in question.  This also proves the claimed frame
bounds.  On the fixed annuli all comparisons remain polynomial after
the exponential changes.  Curvature and any fixed number of its
covariant derivatives are consequently polynomially bounded.

We spell out why local chart control gives a polynomial global
cover.  Shrink the charts uniformly by powers of $\eps$ so that
every point has a chart containing an intrinsic ball of radius
$\rho_\eps\ge c_1\eps^{a_1}$.  In that chart both the coordinate
domain radius and the upper and lower metric eigenvalue bounds are
controlled by powers of $\eps$.  A coordinate ball of polynomial
radius lies in a fixed smaller intrinsic ball, and its volume density
has a polynomial lower bound.  It follows that every intrinsic ball
of radius, say, $\rho_\eps/8$ has volume at least
$c_2\eps^{a_2}$, after a further common polynomial shrinkage if needed.
A maximal disjoint collection of these balls has at most
$(2c_2\eps^{a_2})^{-1}$ members because the total volume is $1/2$.
Their enlarged balls cover $X$ and still fit in controlled charts.
Coordinate cutoff functions and a normalized partition of unity have
polynomial derivative bounds.  Even the crude overlap bound by the
number of charts is polynomial.  This proves (i) with the asserted
global control.

Euclidean Sobolev estimates on this cover, followed by the polynomial
coordinate and covariant-derivative comparisons, give polynomial
intrinsic Sobolev constants.  Since the dimension is four and
$k\ge6$, derivatives of order at most $\lfloor k/2\rfloor$ are bounded
in $L^\infty$ by the $H^k$ norm with such a constant.  Apply the
intrinsic product rule to each derivative of $h\wedge\widetilde h$,
placing the factor with fewer derivatives in $L^\infty$ and the
other in $L^2$.  This proves \Cref{eq:auxiliary-product}.

Before normalization the residuals vanish on the cores and outside
the gluing annuli.  On the annuli their fixed-coordinate derivatives
are exponentially small; all conversions just described cost only
powers of $\eps^{-1}$.  The constant corrections in
\Cref{eq:auxiliary-normalization} introduce exponentially small
coefficients multiplying forms with polynomial intrinsic bounds.
The fixed total volume then gives
\Cref{eq:auxiliary-residual-estimate} and proves (iv).
\end{proof}

\subsection{An inverse on exact two-forms}

The three moment equations now allow us to invert the linearized wedge
equations on exact two-forms.  Stokes' theorem will control their full
$L^2$ norm by their self-dual part.  Estimating the exact two-form itself
in this way avoids a spectral estimate for its one-form primitive in
the collapsing metrics.  The exact-form elliptic framework and the
Stokes identity occur in Donaldson's study of closed two-forms
\cite[Section~2, Proposition~1, and Section~4, Lemma~1]{Donaldson2006};
the estimates below track the constants through the collapsing family.

\begin{lemma}[Exact-form inverse]
\label{lem:auxiliary-inverse}
For fixed $\eps$, let $f=(f_1,f_2,f_3)$ be top forms with
$\int_X f_i=0$.  There is a unique exact two-form
$h=\mathcal L_\eps f$ satisfying
\begin{equation}
  A\wedge h=f_1,\qquad C\wedge h=f_2,\qquad
  D^0_\eps\wedge h=f_3.
  \label{eq:auxiliary-linear-system}
\end{equation}
For each fixed integer $k\ge0$, this inverse extends to the corresponding
Sobolev spaces and satisfies
\begin{equation}
  \|\mathcal L_\eps f\|_{H^k}
     \le P_\eps\|f\|_{H^k},
  \qquad P_\eps\le C_k\eps^{-p_k}.
  \label{eq:auxiliary-inverse-bound}
\end{equation}
No lower bound for the first positive eigenvalue on one-forms is
required.
\end{lemma}

\begin{proof}
Put $E_1=A$, $E_2=C$, and $E_3=D^0_\eps$, and let $M$ be their
pointwise inner-product Gram matrix for $g_\eps$.  The equations
uniquely specify the self-dual part $\phi=h^+$: if
$f_i=t_i\nu$, then
\begin{equation}
  \phi=\sum_{j=1}^3(M^{-1}t)_jE_j.
  \label{eq:auxiliary-selfdual-data}
\end{equation}
Moreover
\begin{equation}
  \langle\phi,E_i\rangle_{L^2}=\int_X f_i=0.
  \label{eq:auxiliary-cokernel-moments}
\end{equation}
Each $E_i$ is closed and self-dual, hence harmonic.
\Cref{prop:topology} gives $b^+(X)=3$, so these three linearly
independent forms are a basis of the harmonic self-dual forms.
Thus \Cref{eq:auxiliary-cokernel-moments} is exactly the solvability
condition for $\dd^+a=\phi$.

One may see existence directly from the Hodge Green operator $G$ for
$g_\eps$ \cite[Equations~(2.1)--(2.2)]{TaylorHodge2010}.
Since $\phi$ is orthogonal to the harmonic self-dual
forms and the Hodge Laplacian commutes with the star, $G\phi$ is
self-dual and $\Delta G\phi=\phi$.  For a self-dual two-form $\psi$,
$2(\dd\dd^*\psi)^+=\Delta\psi$.  Therefore
\begin{equation}
  h=2\dd\dd^*G\phi
  \label{eq:auxiliary-green-formula}
\end{equation}
is exact and has self-dual part $\phi$.  This formula proves
existence for each fixed metric; we will not estimate the Green
operator itself.

If $h$ is exact, Stokes' theorem gives
\begin{equation}
  0=\int_Xh\wedge h
    =\|h^+\|_2^2-\|h^-\|_2^2,
  \qquad
  \|h\|_2=\sqrt2\,\|h^+\|_2.
  \label{eq:auxiliary-exact-ltwo}
\end{equation}
In particular an exact anti-self-dual two-form is zero, proving
uniqueness.  This identity is also the required $L^2$ estimate,
independent of any positive spectral gap.

For higher derivatives use
\begin{equation}
  \dd h=0,\qquad \dd(*h)=2\dd\phi,\qquad
  \dd^*h=-2*\dd\phi.
  \label{eq:auxiliary-divergence-identities}
\end{equation}
The integrated Weitzenb\"ock identity for two-forms gives
\[
  \|\nabla h\|_2^2
  \le 4\|\dd\phi\|_2^2
      +C\|\operatorname{Rm}\|_\infty\|h\|_2^2.
\]
To iterate, apply the same identity to the tensor-valued form
$\nabla^j h$.  Commuting covariant differentiation with alternating
derivative and divergence produces only curvature derivatives
multiplied by lower derivatives of $h$.  Using
\Cref{eq:auxiliary-divergence-identities}, this yields, for fixed $j$,
\[
  \|\nabla^{j+1}h\|_2
  \le C_j\|\phi\|_{H^{j+1}}
      +K_{j,\eps}\|h\|_{H^j},
\]
where $K_{j,\eps}$ is a polynomial in finitely many suprema of
curvature and its derivatives.  Those suprema are polynomially
bounded by \Cref{lem:auxiliary-geometry}.  Induction starting with
\Cref{eq:auxiliary-exact-ltwo} thus bounds $\|h\|_{H^k}$ by a
polynomial multiple of $\|\phi\|_{H^k}$, without losing derivatives.
Finally \Cref{eq:auxiliary-selfdual-data} and the frame and
inverse-Gram bounds give the same estimate in terms of $f$.
Smooth approximation, or the fixed-metric Green formula, gives the
Sobolev version.  Uniqueness is always asserted for $h$, rather than
for its one-form primitive.
\end{proof}

\subsection{The nonlinear correction}

\begin{proof}[Proof of \Cref{thm:auxiliary-solution}]
Fix once and for all an integer $k\ge6$.  Every constant and exponent
in the following argument uses only finitely many derivatives at
this order.  Let $\mathcal E^k_\eps$ be the closed subspace of
$H^k$ two-forms consisting of exact forms.  Equivalently, its elements
are distributionally closed and have zero harmonic projection; they
possess $H^{k+1}$ primitives for this fixed metric.

Writing $D_\eps=D^0_\eps+h$, the desired equations become the
fixed-point equation
\begin{equation}
  h=\mathcal T_\eps(h):=
   \mathcal L_\eps
    \left(-A\wedge D^0_\eps,\,
          -C\wedge D^0_\eps,\,
          \frac{A^2-(D^0_\eps)^2-h\wedge h}{2}\right).
  \label{eq:auxiliary-contraction-map}
\end{equation}
This map is defined on all of $\mathcal E^k_\eps$.  Indeed, the
first two data have integral zero by
\Cref{eq:auxiliary-exact-moments}, as does the free term in the
third slot.  For each exact iterate, including Sobolev iterates,
$\int_Xh\wedge h=0$ by approximation and Stokes' theorem.
Thus every iteration satisfies precisely the three mean conditions
of \Cref{lem:auxiliary-inverse}.

Let $\delta_\eps=\|\mathcal R_\eps\|_{H^k}$, let $P_\eps$ be an
inverse bound from \Cref{eq:auxiliary-inverse-bound}, and let
$Q_\eps$ be a multiplication bound from
\Cref{eq:auxiliary-product}.  Increasing them if necessary, write
\begin{equation}
  P_\eps\le C\eps^{-p},\qquad
  Q_\eps\le C\eps^{-q},\qquad
  \delta_\eps\le C\eps^{-s}e^{-c/\eps}.
  \label{eq:auxiliary-polynomial-losses}
\end{equation}
On the ball of radius $R_\eps=2P_\eps\delta_\eps$ in
$\mathcal E^k_\eps$,
\begin{align}
  \|\mathcal T_\eps(h)\|_{H^k}
  &\le P_\eps\delta_\eps+
          \tfrac12P_\eps Q_\eps R_\eps^2,
     \label{eq:auxiliary-ball-bound}\\
  \|\mathcal T_\eps(h)-\mathcal T_\eps(\widetilde h)\|_{H^k}
  &\le P_\eps Q_\eps R_\eps\,
           \|h-\widetilde h\|_{H^k}.
     \label{eq:auxiliary-lipschitz-bound}
\end{align}
Choose $\eps$ so small that
\begin{equation}
  2P_\eps^2Q_\eps\delta_\eps<\frac12.
  \label{eq:auxiliary-smallness}
\end{equation}
This is possible because its left side is at most
$C\eps^{-(2p+q+s)}e^{-c/\eps}$, which tends to zero.
The ball is invariant and the Lipschitz constant is less than
$1/2$.  The contraction theorem gives an exact solution with
\begin{equation}
  \|h\|_{H^k}\le 2P_\eps\delta_\eps.
  \label{eq:auxiliary-correction-size}
\end{equation}
If the residual is zero, the same conclusion follows directly by
taking $h=0$.  Polynomial losses can always be absorbed by replacing
$c$ with any fixed smaller positive exponent.  Thus the correction
is exponentially small also in $C^0$, by the polynomial Sobolev
embedding bound.  The uniformly positive Gram matrix of the
approximate triple stays positive, with the same choice of sign.
\Cref{eq:auxiliary-contraction-map} gives
\Cref{eq:auxiliary-exact-equations} exactly.

We finish by proving spatial smoothness; no simultaneous estimate
over infinitely many derivative orders is needed.  For the now fixed
$\eps$, choose a Coulomb primitive $a\in H^{k+1}$ of $h$, so
$\dd a=h$ and $\dd^*a=0$.  The three equations for
$D^0_\eps+\dd a$, together with this gauge, form a first-order
nonlinear system for the four components of $a$.
At the solution the principal symbol of its linearization sends a covector-valued
unknown $b$ at a nonzero covector $\xi$ to the three wedge slots
of $\xi\wedge b$ against $A,C,D_\eps$, together with
$\iota_{\xi^\sharp}b$.
If the three wedge slots vanish, $\xi\wedge b$ belongs to the
negative-definite wedge-orthogonal complement of the positive
three-plane.  But $(\xi\wedge b)^2=0$, so $\xi\wedge b=0$.
It follows that $b$ is a multiple of $\xi$, and the gauge slot
then forces $b=0$.  The symbol is square, hence invertible.

Since $k\geq6$ in real dimension four, the initial Sobolev order gives
$a\in C^{4,\alpha}$ and at least $C^{3,\alpha}$ coefficients
for this elliptic linearization, for some $0<\alpha<1$.
Differentiate the system in a smooth coordinate chart.  The
derivative of $a$ satisfies the linearized elliptic system, with
right-hand side involving only smooth background derivatives and
already controlled first derivatives of $a$.  These coefficients and
the right-hand side are $C^{3,\alpha}$, so local first-order elliptic
estimates with finitely differentiable coefficients
\cite[Section~4, Theorem~4.3]{TaylorNonlinearEllipticNotes}
give $\partial a\in C^{4,\alpha}$ and hence
$a\in C^{5,\alpha}$.  Repeating this argument gives smoothness.
All constants in this bootstrap may depend on the fixed $\eps$.
The resulting $D_\eps$ is therefore smooth and proves the theorem.
\end{proof}

\subsection{The auxiliary class and a uniform chamber estimate}

We now prove \Cref{prop:auxiliary-class}.  The prepared fibration lets
us compute the class modulo $\R F$.  The resulting estimate will apply
to every later definite pair with the same periods, even when the fibers
are no longer complex lines for that pair.

\begin{proof}[Proof of \Cref{prop:auxiliary-class}]
We identify integral curve classes with their Poincar\'e duals.
Choose disjoint closed nodal neighborhoods $N_i$ whose boundaries
lie outside all gluing cutoffs, and put
$U=X\setminus\bigcup_i\operatorname{int}N_i$.  The restriction of
$E_\eps$ to $U$ is $D^{\mathrm{sf}}_\eps$.  If there is at least
one node, $U$ fibers over a compact surface with nonempty boundary,
whose second real cohomology is zero.  Thus the basic term in
\Cref{eq:auxiliary-semiflat} is exact on $U$, and
\begin{equation}
  [D^0_\eps]|_U
  =r_\eps\bigl(\eps[\Psi]|_U
         -a_\eps[A]|_U-b_\eps[C]|_U\bigr).
  \label{eq:auxiliary-restriction-class}
\end{equation}
In particular its restriction divided by $\eps$ is bounded.

When there is a node, the kernel of the restriction map
$H^2(X;\R)\longrightarrow H^2(U;\R)$ is exactly $\R F$.
To verify this assertion, excision and
Poincar\'e--Lefschetz duality identify
\[
  H^2(X,U;\R)
    \simeq\bigoplus_i H^2(N_i,\partial N_i;\R)
    \simeq\bigoplus_i H_2(N_i;\R).
\]
A one-node neighborhood has the homotopy type of a torus with its
primitive vanishing cycle killed, and its second homology is generated
by the fiber.  Under the map to $H^2(X;\R)$ these generators give
the Poincar\'e dual fiber classes, all equal to the nonzero class $F$.
The long exact sequence of the pair gives the assertion.
Hence restriction induces an injective map from
$H^2(X;\R)/\R F$ onto its image.  Its inverse on that
finite-dimensional image is bounded.  \Cref{eq:auxiliary-restriction-class} proves that the class
$[D^0_\eps]$ modulo $\R F$ is $O(\eps)$.

The nonlinear correction is exact, so
$[D_\eps]=[D^0_\eps]\in V$.  Choose a fixed linear complement of
$\R F$ in $V$ and use it to lift the quotient class divided by
$\eps$ to a bounded vector $u_\eps$.  This gives
$[D_\eps]=\lambda_\eps F+\eps u_\eps$.
Choose any regular fiber outside the gluing region.  The original
glued form has integral $\eps$ there; the subtracted forms $A,C$
have zero fiber integrals, and the correction is exact.  Therefore
\begin{equation}
  [D_\eps]F=r_\eps\eps,\qquad u_\eps F=r_\eps=1+O(e^{-c/\eps}).
  \label{eq:auxiliary-fiber-area}
\end{equation}
The square normalization and $F^2=0$ now give
\begin{equation}
  \lambda_\eps
    =\frac{1-\eps^2u_\eps^2}{2\eps(u_\eps F)}.
  \label{eq:auxiliary-lambda}
\end{equation}
Boundedness of $u_\eps$ and
\Cref{eq:auxiliary-fiber-area} prove
\Cref{eq:auxiliary-class-expansion}.

If there are no nodal fibers, $\Psi$ is defined on all of $X$ and
the basic term in \Cref{eq:auxiliary-semiflat} is a multiple of
$F$ in cohomology.  The form $D^{\mathrm{sf}}_\eps$ already solves
the equations, so one takes it for $D_\eps$; the same conclusions
follow with $u_\eps=[\Psi]$, $r_\eps=1$, and the coefficient
computed by \Cref{eq:auxiliary-lambda}.

It remains to establish the uniform implication.  By
\Cref{cor:curve-signs}, every curve class under consideration with
$[D_\eps]d>0$ satisfies $d^2\ge-2$ and the following exhaustive
alternatives: either $m=Fd\ne0$ has the same sign as $[B]d$,
or $d=cF$ for a nonzero real number $c$.  The latter alternative
has no sign difficulty, since both $[D_\eps]F$ and $[B]F$ are
positive.

Choose a null vector $G_1\in V$ with $FG_1=1$.  Such a choice is
obtained from any vector pairing to one with $F$ by subtracting
half its square times $F$.  The orthogonal complement of
$\operatorname{span}(F,G_1)$ in the Lorentzian space $V$ is
negative definite.  Decompose
\[
  d=mG_1+nF+z,\qquad
  u_\eps=\alpha_\eps G_1+\beta_\eps F+w_\eps,
\]
where $z,w_\eps$ lie in that complement, and use the Euclidean norm
$|z|^2=-z^2$ there.  We have
$\alpha_\eps=u_\eps F\ge a_0$ and uniform bounds for
$\alpha_\eps,\beta_\eps,w_\eps$.

Suppose that $m<0$.  Integrality of $F$ and $d$ means
$m=-a$ for an integer $a\ge1$.  The square bound gives
\begin{equation}
  -2an-|z|^2=d^2\ge-2,
  \qquad
  n\le\frac{2-|z|^2}{2a}.
  \label{eq:auxiliary-negative-intersection}
\end{equation}
Complete the square in the negative-definite component:
\begin{align}
  u_\eps d
   &\le\frac{\alpha_\eps}{a}
      -\frac{\alpha_\eps|z|^2}{2a}
      +|w_\eps||z|+|\beta_\eps|a \notag\\
   &\le\frac{\alpha_\eps}{a}
      +a\left(\frac{|w_\eps|^2}{2\alpha_\eps}
                    +|\beta_\eps|\right)
    \le Ka.
  \label{eq:auxiliary-completed-square}
\end{align}
The constant $K$ is independent of $a,z,d$, and of the future
almost-complex structure.  It follows that
\begin{equation}
  [D_\eps]d
   =-\lambda_\eps a+\eps u_\eps d
   \le a\left(-\frac{c_0}{\eps}+K\eps\right)<0
  \label{eq:auxiliary-chamber-exclusion}
\end{equation}
whenever $\eps^2<c_0/K$ (with no restriction from this inequality
if $K=0$).  This contradicts $[D_\eps]d>0$ and excludes every
negative integer $m$ simultaneously.  The remaining alternative
$m>0$ gives $[B]d>0$, proving
\Cref{eq:auxiliary-uniform-transfer}.

Finally, both $[D_\eps]$ and $[B]$ have positive square and
pair positively with the same nonzero null vector $F$.  In a
Lorentzian space the two components of the positive cone are
distinguished by the sign of pairing with a fixed nonzero null
vector: no positive vector is orthogonal to it.  Thus the two
classes are in the same component.  All restrictions on $\eps$
in this proof depend only on fixed construction data and the
original periods.  They can therefore be imposed before choosing
any later pair deformation.
\end{proof}
